\begin{abstract}
Generative art has long borrowed emergence from abstract rule systems: cellular automata, flocking, or the simple forms and behaviours of Casey Reas's \emph{Process Compendium}. Their rules are elegant, but arbitrary: they are chosen for what they produce, not for what they are about. This paper presents \emph{emergence-study}, an artwork and software system whose rules are instead drawn from the rhythms of everyday life. Beings in a simulated world are born, age, keep daily routines between places, stay, have children and die; how busy, how energetic and how long-lived they are depends on their age, with noise. The system itself is never shown. What is shown are \emph{software-interpretations}: each begins as a thought, written in English from the artist's lived experience, and is realised as an algorithm that reads selected data from the system and draws it. Every interpretation is documented in three parts, \emph{thought}, \emph{expression} and \emph{parameters}, a structure that extends the instruction-based practices of Sol LeWitt and Yoko Ono to a running simulation. I describe the base system in full, the interpretation method, and six interpretations made with it. I also report how the base system was revised after measurement, and show that the rules of the ground, not only the choices of the reader, decide what an interpretation can reveal. The work argues that a system built from everyday life lets one set of rules hold many readings that are at once personal and grounded, and it offers the three-part interpretation as a reusable, legible format for artists working with simulations.
\end{abstract}

\noindent\textbf{Keywords:} emergence, generative art, instruction art, agent-based simulation, software art, everyday life, process art, data representation.
